\section{Сравнение примененных методов верификации}

На основе полученных данных проведем комплексное сравнение рассмотренных методов верификации, особое внимание уделяя различиям между симуляцией и формальной верификацией. По умолчанию выводы делаются для устройств, для которых применима формальная верификация, если не оговорено иное. 

\subsection{Прогресс верификации}

\tab Для оценки прогресса верификации выберем два основных параметра: 

\begin{enumerate}
    \item Время, затрачиваемое от начала разработки устройства до обнаружения ошибки.
    
    \item Время, затрачиваемое на полную верификацию устройства. 
\end{enumerate}

Разумеется, точные цифры зависят от многих факторов, таких как опыт верификатора, сложность блока, качество спецификации, вычислительные ресурсы и др. Однако на основе практического применения и общих соображений мы можем качественно сравнить два метода. 

Если целью является как можно быстрее обнаружить ошибку, то, как правило, лучше использовать формальную верификацию. Приведем несколько аргументов: 

\begin{itemize}
    \item В отличие от симуляции, разработку формального проекта можно начинать даже до завершения модуля RTL-инженером, ведь набор формальных свойств --- это и есть спецификация, просто записанная на формальном языке. При этом утверждения легко можно включать или отключать, чтобы сосредоточиться на одной части дизайна. 
    \item Судить о корректности работы блока при использовании симуляции можно лишь при почти завершившейся разработке тестбенча и последовательностей. Формальная верификация обладает инкрементальностью, поэтому результаты можно увидеть почти мгновенно. Этому также способствует возможность переиспользовать свойства для разных блоков (например, для valid-then-ready handshake).
\end{itemize}

Для полной верификации устройства ответ разнится в зависимости от того, где мы устанавливаем границу уверенности в том, что устройство работает корректно. Основываясь на практическом опыте, можно сказать, что разработка формальных свойств занимает время, сопоставимое с разработкой верификационного окружения. Тем не менее, симуляция занимает гораздо меньше времени на проведение всех тестов и сбор покрытия (минуты по сравнению с часами или днями). В связи с этим первые предположения о корректном поведении системы у динамических методов появляются заметно раньше. 

В качестве подтверждения сведём в таблицу $ $ \ref{tab:main-table} данные, полученные в ходе применения методов формальной верификации и симуляции:

\begin{table}[h]
	\begin{tabular}{|p{5cm}|| >{\centering\arraybackslash}p{2.75cm}|c|c|c|c|}
		\hline
		Название блока                                                              & D-stage            & C-stage & DD-stage & T-stage & Q-stage \\ \hline
		Количество триггеров                                                        & 374086             & 22457   & 1788     & 79936   & 35906   \\ \hline
		Количество логических вентилей                                              & 1074060            & 142640  & 34599    & 271233  & 195514  \\ \hline
		Время, затраченное на формальную верификацию, минут                         & 4694 (ограниченно) & 1031    & 46       & 3752    & 720     \\ \hline
		Количество транзакций, необходимое для сбора 100\% покрытия & 45 & 150 & 18 & 2000 & 390 \\ \hline 
		Время, затрачиваемое симуляцией для сбора полного кодового покрытия, секунд & 10                 & 12      & 10       & 201     & 19      \\ \hline
	\end{tabular}
	\caption{Время, затраченное на верификацию блоков используя методы формальной верификации и UVM-тестирования.}
	\label{tab:main-table}
\end{table}

Тем не менее, симуляция не способна предоставить исчерпывающий ответ в какие-либо разумные временные рамки. Пусть мы имеем блок сложностью примерно $10^{20}$ (около $2^{66}$-$2^{67}$) состояний и симулятор, который способен симулировать 1 миллиард тактов в секунду, что сильно превышает возможности современных симуляторов. В таком случае на проверку всех возможных состояний нам потребуется около 3 тысяч лет. В то же время формальная верификация предоставила ответ для D-stage, верхнюю границу сложности которого можно оценить в $2^{374086}$ состояний, в течение 3 суток. 

\subsection{Качество верификации}

В целом, для систем среднего размера можно сказать, что ни один из методов не может обеспечить полную верификацию устройства, так как все они упираются в вычислительные ресурсы: проверка всех возможных состояний симуляцией занимает, как показано ранее, миллионы лет, а формальные методы будут требовать слишком много времени и памяти, в связи с чем придется так или иначе упрощать устройство. 

Тем не менее, здесь формальная верификация обладает существенным преимуществом. Под заданными предположениями инструмент проверяет всё пространство возможных состояний системы, включая сложные для обнаружения симуляцией угловые случаи. Иначе говоря, устройство проверяется \textit{исчерпывающе}. 

\subsection{Исправление ошибок}

Часто обнаружить ошибку бывает гораздо легче, чем найти её причину, поэтому стоит обратить внимание на то, в какой форме метод сообщает об ошибке и какие возможности для её исправления ("дебага") он даёт.

Симуляция и в случае успеха, и в случае провала предоставляет фиксированную, очень длинную (тысячи или сотни тысяч тактов) вейвформу с большим числом нерелевантной информации. Это может быть полезно для визуализации поведения блока, однако для обнаружения причины ошибки крайне неудобно. Симуляция сообщает, что ожидаемые и реальные выходные данные не совпадают, но не сообщает, почему. Более того, ошибка может возникать настолько редко, что возможность воспроизвести её при каких-либо других входных данных становится маловероятной. 

С другой стороны, формальная верификация почти всегда предоставляет кратчайший путь к достижению ошибки. Современные инструменты формальной верификации указывают на логику, влияющую на результат (COI, Proof Core), а также позволяют изменить вейвформу для изучения поведения блока. Таким образом, улучшается локальность обнаружения ошибки. 

\subsection{Порог входа}

Порог входа в симуляцию может быть разным в зависимости от подхода: так, например, использовать направленное тестирование гораздо легче, чем ограниченное случайное. Тем не менее, задачу сильно облегчает наличие широко используемой библиотеки UVM, и интуитивность подхода.

У формальной верификации порог входа на порядок выше. В первую очередь это вызвано строгими требованиями к записи спецификации вместе с необходимостью глубокого изучения SVA. От верификатора требуется чёткое понимание того, как именно преобразовать спецификацию в набор свойств, чтобы он был достаточным для проверки модели. При упрощении устройства требуется особая осторожность, чтобы не ограничить пространство состояний слишком сильно и не пропустить ошибку. Также зачастую требуется опыт в оптимизации свойств, которая сама по себе представляет отдельную подзадачу. 

\subsection{Применение}

На основании теоретического исследования и практических результатов можем составить некоторый список практических рекомендаций по применению различных методов верификации. В целом, их можно свести к грамотному распределению ресурсов: формально верифицировать можно и сложные блоки, но на разработку свойств и их доказательство уйдёт очень много времени и вычислительных мощностей, не говоря о том, что свойства будут доказаны лишь для первых нескольких тактов. В то же время, где это возможно, хотелось бы иметь полную уверенность в корректности поведения устройства. 

Поэтому в первую очередь следует проверить сложность модели. Если блок большой (сложность выходит за пределы сотен тысяч триггеров), то нет смысла рассматривать применение формальной верификации и стоит ограничиться симуляцией. Более того, если блок обладает достаточно высоким уровнем абстракции (в контексте данной работы это может быть кэш целиком или одна из его партиций), то следует рассмотреть метод трасс в дополнение к ограниченному случайному тестированию. 

Если блок имеет средний размер (условно до 100000 триггеров) и небольшую глубину, то имеет смысл рассмотреть комбинацию формальной верификации и симуляции. Формальная верификация позволит удостовериться в корректности работы если не всего блока целом, то хотя бы его критически важных частей, где нам важно удостовериться в отсутствии ошибок. Симуляция же является более традиционным методом для таких устройств, и позволяет закрепить результат. 

Для малых блоков (до 10000 триггеров) целесообразно использовать только формальную верификацию. Написание свойств для таких блоков навряд ли будет занимать более нескольких дней, как и их проверка. В то же время симуляция на небольших системах может показывать неудовлетворительные результаты: разработка окружения и сбор полного покрытия будут занимать непропорционально много ресурсов. 

В общем случае, во время разработки модуля, когда верификация выполняется RTL-разработчиком, лучше использовать направленное тестирование и/или разработать тривиальные формальные свойства. Оба этих метода обеспечивают быстрые результаты и позволяют провести поверхностное тестирование. 

Между подтверждением первоначальной функциональности модели и окончательной "заморозкой" (завершением разработки) RTL может пройти очень много времени. В этот период велика вероятность изменений в блоке, его полного переписывания или удаления. Полную формальную верификацию проводить можно, но проверка всех свойств может занимать до одной недели. Симуляция, с другой стороны, быстрее выдаст первые результаты: уже в первый день после внесения изменений можно приблизительно судить о наличии ошибок в блоке.

В конце цикла верификации, по возможности, лучше разработать хотя бы несколько наиболее важных свойств блока, особенно для критических компонентов. Это даст высокую степень уверенности в корректности функционирования модели. 

Для часто используемых блоков и/или блоков небольшого размера (FIFO, буферы, арбитры, счётчики и др.) крайне рекомендуется использовать формальную верификацию как основной метод. Обычно такие блоки плохо покрываются симуляцией, когда они находятся в составе других модулей, и ошибки в них могут очень долго не проявляться в регрессии. Так, например, на ранних стадиях изучения формальной верификации автором была найдена ошибка в буфере, при том что регрессионное тестирование блока, содержащего этот буфер, до этого проходило успешно.
